\chapter{API \& intégrations}
\label{ch:api}

mybibli expose une petite API JSON HTTP sous \texttt{/api/v1/*} qui
permet à des scripts externes et à des outils d'intelligence
artificielle de lire votre catalogue et de modifier la classification
des titres. Cette API est servie par le même serveur web que
l'interface humaine — pas de démon séparé, pas de port supplémentaire
à ouvrir.

Ce chapitre couvre :

\begin{itemize}
  \item la création d'une clé d'API dans le panneau d'administration,
  \item les deux portées (\emph{lecture seule} et \emph{lecture +
        écriture}),
  \item les points d'entrée exposés sous \texttt{/api/v1/*},
  \item trois exemples de bout en bout (\texttt{curl}, Python, boucle
        de classificateur IA).
\end{itemize}

\section{Pourquoi une API HTTP ?}

Le cas d'usage quotidien est de déléguer le travail de classification
à un outil externe — par exemple demander à un LLM de suggérer un
code Dewey, un genre ou un sous-titre propre pour chaque nouveau
livre en anglais. Le LLM parcourt le catalogue avec
\texttt{GET /api/v1/titles}, récupère le détail de chaque entrée via
\texttt{GET /api/v1/titles/\{id\}}, calcule une suggestion, puis
l'écrit avec \texttt{PATCH /api/v1/titles/\{id\}}. Chaque PATCH est
journalisé pour pouvoir revenir en arrière si un modèle a mal
travaillé.

L'API n'est \emph{pas} prévue pour l'import en masse à haut débit —
pour cela, utilisez directement la base SQL lors d'une fenêtre de
maintenance. Elle \emph{est} prévue pour le genre d'écritures lentes,
réfléchies et opinionées qui s'insèrent dans une boucle d'agent LLM.

\section{Créer une clé d'API}
\label{sec:api-keys-create}

Les clés sont gérées dans \textbf{Administration \(\rightarrow\)
Clés d'API}. Seuls les utilisateurs Admin voient cet onglet.

\begin{enumerate}
  \item Cliquez sur \textbf{Administration} dans la barre de
        navigation.
  \item Ouvrez l'onglet \textbf{Clés d'API}.
  \item Saisissez un \textbf{libellé} — un nom que vous seul voyez,
        par ex.\ \emph{« Classificateur IA »} ou \emph{« Script de
        sauvegarde »}. Le libellé sert au journal d'audit ; la clé
        elle-même est aléatoire et indépendante du libellé.
  \item Choisissez une \textbf{portée} :
        \begin{itemize}
          \item \textbf{Lecture seule} — uniquement les points
                d'entrée \texttt{GET}. Idéal pour des tableaux de
                bord en lecture, des outils de miroir ou des
                scripts de curiosité.
          \item \textbf{Lecture + écriture} — autorise en plus
                \texttt{PATCH /api/v1/titles/\{id\}} pour une
                liste blanche restreinte de champs (genre, Dewey,
                description, sous-titre). À utiliser pour des
                agents LLM ou des outils de curation.
        \end{itemize}
  \item Cliquez sur \textbf{Créer la clé}. Une boîte de dialogue
        affiche la valeur complète une seule fois — copiez-la
        maintenant. Une fois fermée, il n'y a aucun moyen de la
        récupérer ; si vous perdez la clé, révoquez-la et créez-en
        une nouvelle.
\end{enumerate}

La clé a la forme \texttt{mybibli\_<portée>\_<aléa40>}, où
\texttt{<portée>} vaut \texttt{ro} ou \texttt{rw}. Seuls le hachage
argon2 et les 12 premiers caractères du texte clair (le préfixe)
sont persistés en base ; le préfixe sert à restreindre les candidats
lors de l'authentification. Le hachage ne peut pas être inversé.

\section{S'authentifier}

Chaque requête \texttt{/api/v1/*} doit transporter l'un des deux
en-têtes :

\begin{itemize}
  \item \texttt{Authorization: Bearer <clé>} — préféré.
  \item \texttt{X-API-Key: <clé>} — solution de repli pratique.
\end{itemize}

Une clé absente ou invalide retourne \texttt{401 Unauthorized} avec
une enveloppe JSON :

\begin{verbatim}
{ "error": "unauthorized",
  "reason": "missing_api_key" }
\end{verbatim}

Une clé en lecture seule sur un point d'entrée d'écriture retourne
\texttt{403 Forbidden} avec la portée manquante explicitement nommée :

\begin{verbatim}
{ "error": "forbidden",
  "reason": "insufficient_scope",
  "scope_required": "write" }
\end{verbatim}

La protection CSRF est volontairement court-circuitée pour
\texttt{/api/*} parce que l'authentification par jeton ne repose
pas sur les cookies, et le risque de rejeu cross-site contre lequel
CSRF protège ne s'applique pas.

\section{Points d'entrée}

\subsection{GET /api/v1/titles}

Liste paginée. Accepte \texttt{?page=N},
\texttt{?q=<recherche>}, \texttt{?genre\_id=N} et
\texttt{?dewey\_prefix=8}.

Retour :

\begin{verbatim}
{ "items": [ { "id": 42, "title": "...",
               "subtitle": null,
               "language": "fr",
               "media_type": "book",
               "isbn": "9782...",
               "publisher": "...",
               "publication_year": 2020,
               "genre_id": 7,
               "genre_name": "Roman",
               "dewey_code": "843.92" },
             ... ],
  "page": 1, "total_pages": 12,
  "total_items": 287 }
\end{verbatim}

\subsection{GET /api/v1/titles/\{id\}}

Détail complet, incluant les contributeurs, les volumes (avec leur
emplacement) et les affectations de série. Le champ
\texttt{version} est le verrou optimiste — conservez-le à côté des
données que votre script va modifier.

\subsection{GET /api/v1/genres, /locations, /series}

Données de référence. À utiliser quand un LLM a besoin de connaître
les genres existants avant de suggérer un \texttt{genre\_id}.

\subsection{GET /api/v1/dewey/\{prefix\}}

Recherche par bucket. Passez un préfixe Dewey partiel (par ex.
\texttt{8} pour la littérature) et récupérez tous les titres dont
le \texttt{dewey\_code} commence par ce préfixe. Utile pour
« trouver tous les livres qui ont déjà un code Dewey en 800 que je
puisse comparer ».

\subsection{PATCH /api/v1/titles/\{id\}}

\textbf{Portée écriture uniquement.} Modifie les quatre champs en
liste blanche :

\begin{itemize}
  \item \texttt{subtitle} — peut être null.
  \item \texttt{description} — peut être null.
  \item \texttt{dewey\_code} — peut être null ; chiffres, points
        et slashes uniquement.
  \item \texttt{genre\_id} — doit référencer un genre existant
        et non supprimé.
\end{itemize}

Le corps JSON suit une sémantique de mise à jour partielle :

\begin{itemize}
  \item Champ \textbf{omis} \(\rightarrow\) pas de changement.
  \item Champ à \texttt{null} \(\rightarrow\) effacement vers NULL
        (pour les trois champs nullable ; \texttt{null} sur
        \texttt{genre\_id} retourne \texttt{400}).
  \item Champ avec une valeur \(\rightarrow\) mise à jour.
\end{itemize}

\texttt{version} est obligatoire. En cas de désaccord la réponse
est \texttt{409 Conflict} et contient la version attendue et la
version reçue — relisez la donnée et réessayez.

Un PATCH réussi écrit une ligne dans \texttt{admin\_audit} avec
l'action \texttt{api\_patch\_title}, l'attribution via l'admin
émetteur (\texttt{user\_id}) et la clé d'API
(\texttt{details.via\_api\_key\_id} +
\texttt{details.via\_api\_key\_label}), et un diff
\texttt{old}/\texttt{new} par champ. Un PATCH sans changement
effectif renvoie \texttt{200} mais NE bump PAS \texttt{version} et
N'ÉCRIT PAS de ligne d'audit.

\section{Exemples}

\subsection{\texttt{curl}}

\begin{verbatim}
KEY="mybibli_ro_..."

# Lister la première page des titres.
curl -H "Authorization: Bearer $KEY" \
     https://library.example.org/api/v1/titles

# Afficher un titre en détail.
curl -H "Authorization: Bearer $KEY" \
     https://library.example.org/api/v1/titles/42

# PATCH un code Dewey (clé write requise).
KEY_WRITE="mybibli_rw_..."
curl -X PATCH \
     -H "Authorization: Bearer $KEY_WRITE" \
     -H "Content-Type: application/json" \
     -d '{"version": 3, "dewey_code": "813.54"}' \
     https://library.example.org/api/v1/titles/42
\end{verbatim}

\subsection{Python}

\begin{verbatim}
import os, requests

BASE = "https://library.example.org/api/v1"
KEY  = os.environ["MYBIBLI_API_KEY"]
H    = {"Authorization": f"Bearer {KEY}"}

# Parcourir toutes les pages.
page = 1
while True:
    r = requests.get(f"{BASE}/titles?page={page}",
                     headers=H, timeout=30)
    r.raise_for_status()
    data = r.json()
    for t in data["items"]:
        print(t["id"], t["title"], t.get("dewey_code"))
    if page >= data["total_pages"]:
        break
    page += 1
\end{verbatim}

\subsection{Une boucle de classificateur LLM}

Le flux d'écriture prévu ressemble à ceci — en pseudo-code :

\begin{verbatim}
for t in walk_titles(KEY_READ):
    if t["dewey_code"]:
        continue                   # déjà classifié
    detail = get(f"/titles/{t['id']}")
    suggestion = ask_llm(detail)   # renvoie {"dewey_code": ...}
    if not suggestion.get("dewey_code"):
        continue
    patch(
        f"/titles/{t['id']}",
        json = {
            "version":    detail["version"],
            "dewey_code": suggestion["dewey_code"],
        },
        key  = KEY_WRITE,
    )
\end{verbatim}

Si deux classificateurs entrent en course sur le même titre, le
perdant reçoit un \texttt{409} et re-lit la donnée ; pas
d'écrasement silencieux. Le journal d'audit conserve qui (clé)
a suggéré quoi — pratique pour annuler un modèle qui aurait eu
une mauvaise journée.

\section{Révoquer une clé}

La révocation se fait en un clic dans le panneau d'administration.
Une clé révoquée \emph{ne peut pas} être restaurée ; l'historique
d'audit des utilisations passées est préservé à côté de la ligne
désormais désactivée. Si vous soupçonnez une fuite (commit dans un
dépôt public, envoi par e-mail par erreur, présence dans un fichier
de log), révoquez immédiatement et créez-en une nouvelle.

\section{Ce que l'API n'expose \emph{pas}}

Par construction, l'API v1 est majoritairement en lecture. Elle
ne permet pas :

\begin{itemize}
  \item de créer ou supprimer des titres, volumes, contributeurs,
        emprunteurs, prêts ou séries ;
  \item de modifier un champ en-dehors de la liste blanche des
        quatre champs du PATCH ;
  \item de retourner des mots de passe, hachages, jetons de
        session ou autres credentials ;
  \item d'accepter des uploads multipart (pas de PUT de jaquette).
\end{itemize}

Ce sont des limites volontaires de la v1 — le cas d'usage qui
justifie l'API est la classification, et cela tient dans la liste
blanche restreinte. Une surface d'écriture plus large changerait la
posture de sécurité (considérations CSRF, limitation de débit pour
les mutations en masse) et est reportée tant qu'un besoin concret
ne s'est pas manifesté.
